\chapter{The Land Surface}\label{chap:ch12_land}

% --- local symbols for this chapter only (kept minimal) ---
\providecommand{\Tsfc}{T_{\mathrm s}}        % land surface (skin) temperature
\providecommand{\Tair}{T_{\mathrm a}}        % near-surface air temperature
\providecommand{\Wsoil}{W}                    % soil-moisture column depth [m]
\providecommand{\Wmax}{W_{\max}}              % bucket capacity [m]
\providecommand{\Rnet}{R_{\mathrm n}}         % net radiation [W m^{-2}]
\providecommand{\Hsen}{H}                      % sensible heat flux
\providecommand{\LE}{L_v E}                    % latent heat flux
\providecommand{\lai}{L}                       % leaf-area index
\providecommand{\Topt}{T_{\mathrm{opt}}}      % optimal growth temperature

The land surface in \model{} is deliberately the simplest of the three media.
Where the ocean (\ref{chap:ch04_ocean_equations}) carries a full
three-dimensional primitive-equation state and the atmosphere
(\ref{chap:ch09_atmosphere_multilevel}) carries a spectral dynamical core, land is a
\emph{single-layer slab}: one prognostic skin temperature, one column of soil
water held in a leaky bucket, a prognostic leaf-area index, a soil-carbon pool,
and a snow reservoir. There is no horizontal transport within the land --- each
grid column is an independent ``tile'' --- and the only land-to-land coupling is
through the shared atmosphere above it. This is the classic \emph{bucket model}
of Manabe, dressed with a dynamic-vegetation albedo and a one-bucket snow scheme.
Its purpose is to close the surface energy and water budgets seen by the
atmosphere, not to be a land-carbon-cycle laboratory; we are explicit below about
where the simplifications bite.

The implementation is a single \jax{}, fully differentiable time step,
\codref{chronos\_esm/land/driver.py}{\code{step\_land}}, with the vegetation
sub-step in \codref{chronos\_esm/land/vegetation.py}{\code{step\_vegetation}}.

\section{State variables and timescales}

The land state is the tuple
\begin{equation}
  \mathcal{L} = \bigl(\,\Tsfc,\ \Wsoil,\ \lai,\ C_{\mathrm s},\ h_{\mathrm{sn}}\,\bigr),
\end{equation}
namely the surface temperature $\Tsfc$ [K], the soil-moisture depth $\Wsoil$ [m]
(bounded by the bucket capacity $\Wmax = 0.15\ \mathrm{m}$), the leaf-area index
$\lai$ [$\mathrm{m^2\,m^{-2}}$], a soil-carbon density $C_{\mathrm s}$
[$\mathrm{kg\,C\,m^{-2}}$], and a snow water-equivalent depth $h_{\mathrm{sn}}$
[m]. The component is integrated with the \emph{coupling interval} as its time
step (the daily coupling cadence, passed as \code{dt}), not the atmospheric
dynamics step --- a deliberate choice that exploits the long thermal and
hydrological memory of the slab.

\begin{remark}
The slab has a single thermal layer of depth $d=1\ \mathrm{m}$, soil density
$\rho_{\mathrm s}=1600\ \mathrm{kg\,m^{-3}}$ and heat capacity
$c_{\mathrm s}=800\ \mathrm{J\,kg^{-1}\,K^{-1}}$, giving a column heat capacity
$\mathcal{C}_0 = \rho_{\mathrm s} c_{\mathrm s} d \approx
1.28\times10^{6}\ \mathrm{J\,m^{-2}\,K^{-1}}$. The thermal depth was raised from
typical $\sim\!0.1$~m bucket values to 1~m \emph{for numerical stability of the
semi-implicit step}, which makes the diurnal response unphysically sluggish ---
acceptable here because the model is run at daily coupling and validated against
monthly/annual climatology, not the diurnal cycle.
\end{remark}

\section{Surface energy balance}

\subsection{Net radiation}

Land absorbs shortwave according to its albedo and emits longwave as a grey
body. With downwelling shortwave $S^\downarrow$ and longwave $L^\downarrow$
supplied by the atmosphere,
\begin{align}
  S_{\mathrm{net}} &= (1-\alpha)\,S^\downarrow, \label{eq:swnet}\\
  L_{\mathrm{net}} &= L^\downarrow - \varepsilon\,\sigma\,\Tsfc^4,
  \label{eq:lwnet}\\
  \Rnet &= S_{\mathrm{net}} + L_{\mathrm{net}},
\end{align}
where $\sigma$ is the Stefan--Boltzmann constant, $\varepsilon$ the land
emissivity, and $\alpha$ the (dynamic) surface albedo of
Section~\ref{sec:albedo}. See
\codref{chronos\_esm/land/driver.py}{net-radiation block in \code{step\_land}}.

\subsection{Turbulent fluxes}

Sensible and latent heat are bulk-aerodynamic. With air density
$\rho_a=1.225\ \mathrm{kg\,m^{-3}}$, $c_p = 1004\ \mathrm{J\,kg^{-1}\,K^{-1}}$,
wind speed $U$, and a drag coefficient $C_h=C_e$ supplied by the coupler (which
folds in the surface-roughness logic),
\begin{equation}
  \Hsen = \rho_a\,c_p\,C_h\,U\,(\Tsfc - \Tair).
  \label{eq:sensible}
\end{equation}
The latent flux uses a Monteith-style saturation deficit. Saturation specific
humidity is built from the August--Roche--Magnus vapour-pressure formula at the
surface pressure $p_0=101325\ \mathrm{Pa}$,
\begin{align}
  e_s(\Tsfc) &= 611\,\exp\!\Bigl[\frac{17.67\,(\Tsfc-273.15)}
                                       {(\Tsfc-273.15)+243.5}\Bigr], &
  q_{\mathrm{sat}} &= 0.622\,\frac{e_s}{p_0}.
  \label{eq:qsat}
\end{align}
The actual evaporation is the \emph{potential} evaporation throttled by a soil-
and vegetation-dependent moisture availability $\beta$:
\begin{align}
  E_{\mathrm{pot}} &= \max\!\bigl[\rho_a\,C_e\,U\,(q_{\mathrm{sat}}-q_a),\ 0\bigr],
  & E &= \beta\,E_{\mathrm{pot}}, & \LE &= L_v\,E,
  \label{eq:evap}
\end{align}
with the latent heat of vapourisation $L_v = 2.5\times10^{6}\
\mathrm{J\,kg^{-1}}$ (sublimation over snow is treated with the same $L_v$, a
known $\sim\!13\%$ underestimate of the true sublimation enthalpy). The
$\max(\cdot,0)$ forbids dew/condensation. The availability factor is the
\emph{bucket beta}, enhanced by vegetation cover and forced to unity over snow:
\begin{equation}
  \beta = \mathrm{clip}\!\Bigl[\underbrace{\mathrm{clip}\bigl(\tfrac{\Wsoil}{\Wmax},0,1\bigr)}_{\beta_{\mathrm{bucket}}}\,
          (1+f_{\mathrm{veg}}),\ 0,\ 1\Bigr],\qquad
  \beta \leftarrow (1-c_{\mathrm{sn}})\,\beta + c_{\mathrm{sn}}\cdot 1,
  \label{eq:beta}
\end{equation}
where $f_{\mathrm{veg}}$ is the vegetated fraction (Section~\ref{sec:albedo})
and $c_{\mathrm{sn}}$ the snow-cover fraction. The factor $(1+f_{\mathrm{veg}})$
crudely represents transpiration drawing water faster than bare-soil
evaporation, but because the result is clipped to $[0,1]$ it saturates: a
fully-vegetated, half-full bucket already evaporates at the potential rate. See
\codref{chronos\_esm/land/driver.py}{turbulent-flux and $\beta$ block}.

\subsection{Ground heat flux and the semi-implicit temperature update}

The energy reaching the slab is the residual
\begin{equation}
  G = \Rnet - \Hsen - \LE,
\end{equation}
and the slab obeys $\mathcal{C}\,\dt{\Tsfc} = G$, with the total heat capacity
augmented by a snow term,
$\mathcal{C} = \mathcal{C}_0 + \rho_{\mathrm{sn}} c_{\mathrm{sn,heat}}
h_{\mathrm{sn}}$ ($\rho_{\mathrm{sn}}=300$, $c_{\mathrm{sn,heat}}=2000$). Because
$G$ depends strongly and nonlinearly on $\Tsfc$ (through $\Tsfc^4$, $\Hsen$ and
$q_{\mathrm{sat}}$), an explicit Euler step would need a tiny $\Delta t$. The
code instead linearises $G$ about the current temperature and takes one
\emph{semi-implicit} (backward-Euler-in-the-linearised-source) step. Writing
$G(\Tsfc+\Delta\Tsfc)\approx G + G'\,\Delta\Tsfc$ with
\begin{equation}
  G' = \pd{G}{\Tsfc} =
       \underbrace{-4\varepsilon\sigma\Tsfc^3}_{\text{longwave}}
       \;\underbrace{-\,\rho_a c_p C_h U}_{\text{sensible}}
       \;\underbrace{-\,\rho_a L_v C_e U\,\beta\,\pd{q_{\mathrm{sat}}}{\Tsfc}}_{\text{latent}},
  \qquad
  \pd{q_{\mathrm{sat}}}{\Tsfc} = q_{\mathrm{sat}}\,\frac{17.67\cdot 243.5}{(\Tsfc-273.15+243.5)^2},
  \label{eq:dgdt}
\end{equation}
the implicit balance $\mathcal{C}\,\Delta\Tsfc/\Delta t = G + G'\Delta\Tsfc$
gives the increment actually coded,
\begin{equation}
  \boxed{\;\Delta\Tsfc = \frac{G}{\;\mathcal{C}/\Delta t - G'\;}\;}
  \label{eq:dtemp}
\end{equation}
Since $G'<0$ the denominator is always positive and larger than the explicit
$\mathcal{C}/\Delta t$, so the step is unconditionally damped --- this is what
permits the long $\Delta t$. The update $\Tsfc \mapsto \Tsfc + \Delta\Tsfc$ is
applied only on land (multiplied by the land mask). See
\codref{chronos\_esm/land/driver.py}{semi-implicit temperature update}.

\section{Snow and the phase-change correction}

Precipitation is partitioned by the \emph{air} temperature: it is snow where
$\Tair<273.15\ \mathrm{K}$ and rain otherwise (a hard threshold, no mixed
phase). After the temperature update, any column that is both warm
($\Tsfc>273.15$) and snow-covered melts. The available melt energy
$(\Tsfc-273.15)\,\mathcal{C}$ is converted to a melt depth through the latent
heat of fusion $L_f = 3.34\times10^{5}\ \mathrm{J\,kg^{-1}}$,
\begin{equation}
  \Delta h_{\mathrm{melt}} =
  \min\!\Bigl[\max\bigl(0,\ (\Tsfc-273.15)\,\mathcal{C}/L_f/\rho_w\bigr),\ h_{\mathrm{sn}}\Bigr],
\end{equation}
and the slab temperature is corrected \emph{downward} by the latent heat
consumed, $\Tsfc \mapsto \Tsfc - \Delta h_{\mathrm{melt}}\,\rho_w L_f/\mathcal{C}$,
which pins a still-snowy column to the freezing point. The snow reservoir
evolves as
\begin{equation}
  h_{\mathrm{sn}}^{\,n+1} = \max\!\bigl[0,\ h_{\mathrm{sn}}^{\,n}
                            + \Delta t\,P_{\mathrm{snow}} - \Delta h_{\mathrm{melt}}\bigr].
\end{equation}
A final safety clamp $\Tsfc\in[180,340]\ \mathrm{K}$ guards the differentiable
step against transient overshoots. See
\codref{chronos\_esm/land/driver.py}{snow/melt block}.

\section{Bucket hydrology and runoff}\label{sec:bucket}

The soil-water column is a leaky bucket. Inflow is rain plus snowmelt (the melt
depth is converted from a per-step amount back to a rate $\Delta
h_{\mathrm{melt}}/\Delta t$); the only sink is evaporation
$E_{\mathrm{soil}}=E/\rho_w$:
\begin{equation}
  \dt{\Wsoil} = \bigl(P_{\mathrm{rain}} + \Delta h_{\mathrm{melt}}/\Delta t\bigr) - E/\rho_w.
  \label{eq:bucket}
\end{equation}
After the explicit update $\Wsoil^\ast = \Wsoil + \Delta t\,\dt{\Wsoil}$ the
bucket is capped at $\Wmax$, and \emph{everything above the rim is runoff}:
\begin{align}
  R &= \max(\Wsoil^\ast - \Wmax,\ 0), &
  \Wsoil^{\,n+1} &= \mathrm{clip}(\Wsoil^\ast,\ 0,\ \Wmax).
  \label{eq:runoff}
\end{align}
This is instantaneous saturation-excess runoff: there is no infiltration delay,
no sub-surface drainage, and no base-flow recession. Dry-end behaviour is
similarly simple --- the floor at $0$ silently discards any negative residual,
a small non-conservation that is negligible at the validated drift levels.
See \codref{chronos\_esm/land/driver.py}{soil-moisture update and runoff}.

\subsection{Routing runoff to the coastal ocean}

Runoff must reach the ocean or the global water budget leaks. The coupler does
not solve a river network; it performs a one-step \emph{nearest-neighbour
spread} from each land cell to its four neighbours, depositing only into ocean
cells. With the overflow expressed as a freshwater flux $r =
R\,\rho_w/\Delta t$ [$\mathrm{kg\,m^{-2}\,s^{-1}}$] masked to land, the spread is
\begin{equation}
  r_{\mathrm{spread}} = \tfrac14\bigl(r_{i-1,j} + r_{i+1,j} + r_{i,j-1} + r_{i,j+1}\bigr),
  \qquad
  r_{\mathrm{ocean}} = r_{\mathrm{spread}}\cdot m_{\mathrm{ocean}},
  \label{eq:route}
\end{equation}
and $r_{\mathrm{ocean}}$ is added to the ocean's freshwater forcing. The
longitude axis is periodic (a true \code{roll}); the latitude axis uses
\emph{zero-filled} shifts so polar land cannot wrap runoff from the south-pole
row into the north-pole row. See
\codref{chronos\_esm/main.py}{runoff routing in the coupled step}.

\begin{remark}
This scheme conserves freshwater \emph{only} for coastal land. Runoff generated
in a deep continental interior --- whose four neighbours are all land --- is
spread one cell, masked to ocean, and the part landing back on land is
discarded. At T31 most large land masses have such interiors, so a fraction of
terrestrial runoff never reaches the ocean. A genuine flow-accumulation router
(downhill drainage to a single coastal mouth) is the natural upgrade; it is not
yet implemented.
\end{remark}

\section{Vegetation and dynamic albedo}\label{sec:albedo}

Leaf-area index is prognostic and follows a logistic growth law with
temperature- and moisture-limited growth competing against senescence:
\begin{equation}
  \dt{\lai} = \underbrace{r\,g\,\lai\Bigl(1-\frac{\lai}{\lai_{\max}}\Bigr)}_{\text{logistic growth}}
            + \underbrace{s\,g}_{\text{seed}}
            - \underbrace{\Bigl(\frac{1}{\tau_d} + \frac{1-\beta_{\mathrm{bucket}}}{\tau_d/2}\Bigr)\lai}_{\text{decay + drought stress}},
  \label{eq:lai}
\end{equation}
with $\lai_{\max}=6$, growth rate $r=1/\tau_g$ ($\tau_g=10$~d), decay timescale
$\tau_d=30$~d, and a tiny seed term ($s=10^{-8}$) so that bare desert can
re-green when conditions improve. The growth potential $g\in[0,1]$ multiplies a
Gaussian temperature window centred on $\Topt=298.15\ \mathrm{K}$ (width
$15\ \mathrm{K}$, hard-cut to zero below $278.15\ \mathrm{K}$) with the moisture
factor $\beta_{\mathrm{bucket}}$:
\begin{equation}
  g(\Tsfc,\beta_{\mathrm{bucket}}) =
  \beta_{\mathrm{bucket}}\,\exp\!\Bigl[-\frac{(\Tsfc-\Topt)^2}{2\cdot 15^2}\Bigr]
  \cdot \mathbb{1}[\Tsfc \ge 278.15].
\end{equation}
The result is clipped to $[\lai_{\min},\lai_{\max}]=[0.1,6]$. Note this growth
sub-step uses the atmospheric step \code{DT\_ATMOS} internally even though the
thermodynamic step uses the coupling interval --- an inconsistency that makes
the vegetation timescales effectively slower than the nominal $\tau_g,\tau_d$;
it is harmless for climate equilibria but worth knowing before tuning. See
\codref{chronos\_esm/land/vegetation.py}{\code{step\_vegetation}}.

LAI controls albedo through a Beer--Lambert canopy-closure factor. The vegetated
fraction and the area-weighted albedo are
\begin{align}
  f_{\mathrm{veg}} &= 1 - e^{-k\lai}, \quad k=0.5, &
  \alpha_{\mathrm{veg}} &= f_{\mathrm{veg}}\,\alpha_{\mathrm{leaf}}
                          + (1-f_{\mathrm{veg}})\,\alpha_{\mathrm{soil}},
  \label{eq:vegalbedo}
\end{align}
with a dark canopy $\alpha_{\mathrm{leaf}}=0.12$ over bright bare soil
$\alpha_{\mathrm{soil}}=0.35$. The snow-cover fraction then over-rides this,
$c_{\mathrm{sn}}=\mathrm{clip}(h_{\mathrm{sn}}/0.05,0,1)$ (full cover at 5~cm),
giving the final albedo fed back to Eq.~\eqref{eq:swnet},
\begin{equation}
  \alpha = (1-c_{\mathrm{sn}})\,\alpha_{\mathrm{veg}} + c_{\mathrm{sn}}\,\alpha_{\mathrm{snow}},
  \qquad \alpha_{\mathrm{snow}}=0.80.
  \label{eq:albedo}
\end{equation}
This closes a genuine \emph{albedo feedback}: cooling promotes snow, snow raises
$\alpha$, which cools further --- the same positive loop that drives high-latitude
sensitivity in full models, here in its barest form. See
\codref{chronos\_esm/land/vegetation.py}{\code{compute\_land\_properties}}.

\begin{example}[Why the surface is so reflective by default]
At $\lai=1$ (the initial sparse state of \code{init\_land\_state}),
$f_{\mathrm{veg}}=1-e^{-0.5}\approx0.39$, so
$\alpha_{\mathrm{veg}}=0.39\cdot0.12+0.61\cdot0.35\approx0.26$. Even mature
forest ($\lai=6$) gives $f_{\mathrm{veg}}\approx0.95$ and
$\alpha_{\mathrm{veg}}\approx0.13$. Bare/desert tiles therefore sit near
$\alpha=0.35$ --- bright, which is realistic for deserts but means any column
that fails to green stays cold, reinforcing the limit on growth.
\end{example}

\section{Soil carbon: a passive diagnostic}

A single-pool soil-carbon budget is integrated for completeness but is not
coupled back to the atmosphere's $\mathrm{CO_2}$ (which is prescribed; see
\ref{chap:ch14_forcing_response}). Gross primary production is light-use-efficiency
limited, autotrophic respiration is half of GPP, and heterotrophic respiration
is $Q_{10}$-temperature-sensitive,
\begin{align}
  \mathrm{GPP} &= \mathrm{LUE}\,S^\downarrow\,\lai\,\beta, &
  R_a &= 0.5\,\mathrm{GPP}, &
  R_h &= k\,C_{\mathrm s}\,2^{(\Tsfc-T_{\mathrm{ref}})/10}\,\beta, \\
  \dt{C_{\mathrm s}} &= (\mathrm{GPP}-R_a) - R_h, &
  \mathrm{NEE} &= R_a + R_h - \mathrm{GPP},
\end{align}
with $\mathrm{LUE}=10^{-9}$, $k=10^{-8}$, $T_{\mathrm{ref}}=288.15\ \mathrm{K}$,
$C_{\mathrm s}\ge 0$. NEE is returned for diagnostics only. Treat this pool as
an indicator of where a carbon cycle \emph{would} respond, not as a calibrated
flux. See \codref{chronos\_esm/land/driver.py}{carbon-cycle block}.

\section{Summary of approximations}

\begin{table}[h]
\centering
\begin{tabular}{@{}ll@{}}
\toprule
Process & Treatment in \model{} \\
\midrule
Heat storage & 1~m slab, semi-implicit Eq.~\eqref{eq:dtemp}; no diurnal cycle \\
Soil water & single leaky bucket, $\Wmax=0.15$~m, instant saturation runoff \\
Runoff routing & 4-neighbour spread to coastal ocean; interiors leak \\
Snow & one reservoir, $\Tair$ phase split, $L_v$ used for sublimation \\
Vegetation & prognostic LAI (logistic), Beer--Lambert albedo \\
Carbon & passive single pool, not fed back to atmosphere \\
Lateral land transport & none (independent tiles) \\
\bottomrule
\end{tabular}
\caption{The land surface as a closure for the atmospheric boundary, not a
stand-alone biosphere model.}
\end{table}

\begin{exercise}[The bucket time constant]
Using Eqs.~\eqref{eq:evap}--\eqref{eq:bucket}, linearise the bucket about a
half-full state and estimate the e-folding drying time of $\Wsoil$ under constant
potential evaporation with no precipitation. Take $E_{\mathrm{pot}}$ for a
$5\ \mathrm{m\,s^{-1}}$ wind, $C_e=1.5\times10^{-3}$, and a $0.5\%$ humidity
deficit. Compare your estimate to $\Wmax/E_{\mathrm{soil}}$ and explain why the
$(1+f_{\mathrm{veg}})$ enhancement in $\beta$ does \emph{not} shorten it once
$\beta$ saturates at 1.
\end{exercise}

\begin{exercise}[Snow--albedo feedback as a laboratory]
Initialise a high-latitude column at $\Tsfc=274\ \mathrm{K}$ with
$h_{\mathrm{sn}}=0$ and step \code{step\_land} forward with a fixed weak
$S^\downarrow$ and a cold $\Tair=270\ \mathrm{K}$. Show numerically that the
column transitions to a snow-covered, high-albedo state and that the transition
is hysteretic: once snowed in (Eq.~\eqref{eq:albedo} pins $\alpha=0.80$ above
5~cm), a larger $S^\downarrow$ is needed to melt out than was needed to snow in.
Relate this to the bistability ideas of \ref{chap:ch17_amoc_tipping}.
\end{exercise}

\begin{exercise}[Closing the freshwater budget]
Build a synthetic continent (a block of land cells with a land interior) and
verify by direct summation over Eq.~\eqref{eq:route} that the freshwater
delivered to the ocean is \emph{less} than the total runoff $R$ generated, by
the fraction of runoff produced in interior cells. Sketch the modification
needed to make \code{step\_land}'s runoff routing globally conservative (hint:
accumulate interior runoff downstream rather than discarding it).
\end{exercise}

\begin{exercise}[Differentiable land calibration]
The entire land step is written in \jax{} and is differentiable. Pick a single
scalar parameter (e.g.\ $\alpha_{\mathrm{soil}}$ or $\Wmax$), define a loss as
the squared error of seasonal-mean $\Tsfc$ against a target, and use
\code{jax.grad} to compute $\partial \mathcal{L}/\partial \theta$. Comment on
which clamps (the $\Tsfc\in[180,340]$ guard, the $\max(\cdot,0)$ on
evaporation/runoff) produce zero-gradient regions, and how that constrains
gradient-based tuning.
\end{exercise}
